\section{Acerca de este documento}

\subsection{Alcance y propósito}

Este documento presenta la prueba de concepto (PoC) de la cerca virtual del collar DR-VF-01.
Describe lo que hace el firmware, cómo se verificó y qué resultados se midieron, y declara con
la misma claridad lo que la PoC \emph{no} demuestra. Es un informe de ingeniería sobre una parte
del sistema; el proyecto completo, con su arquitectura, hardware, mecánica y cumplimiento, está
en la \emph{Memoria Técnica} (\texttt{docs/DR-VF-01\_Memoria\_Tecnica.pdf}), a la que este informe remite
y que no repite.

La PoC corre en una placa de desarrollo comercial, no en la PCB del producto. Todas las cifras
de este informe se midieron en esa placa y no se pueden trasladar sin más a la PCB final.
Ver el apartado de limitaciones.

\subsection{Público}

Está escrito para un ingeniero que quiera entender el diseño, reproducir las pruebas o partir
de este firmware en otra placa, y para quien evalúe el proyecto.

\subsection{Plataforma y versiones}

\begin{table}[H]
\centering
\small
\begin{tabularx}{\textwidth}{|l X|}
\hline
\tablehead{\textbf{Elemento} & \textbf{Valor}} \\ \hline
Placa & \dBoard \\ \hline
Procesador de la aplicación y GNSS & \dSoc{} (LTE-M, NB-IoT y GNSS) \\ \hline
SDK & \dNcs{} con \dZephyr \\ \hline
Firmware del módem & \dModemFw \\ \hline
Firmware medido & commit \texttt{\dFirmwareCommit} del repositorio \\ \hline
Compilación & \texttt{firmware/build.sh} (ver el apartado de puesta en marcha) \\ \hline
\end{tabularx}
\caption{Plataforma y versiones de la PoC.}
\end{table}

\subsection{Historial de revisiones}

\begin{table}[H]
\centering
\small
\begin{tabularx}{\textwidth}{|c l X|}
\hline
\tablehead{\textbf{Rev.} & \textbf{Estado} & \textbf{Contenido}} \\ \hline
1.0 & Borrador & Firmware de la PoC: cerca cuadrada de 60~m creada con el primer fix GNSS, y modo demo con botón. \\ \hline
1.1 & Planeada & Polígono definido por vértices del usuario (cerca de una finca). Se añadirá como una sección nueva y filas nuevas en las tablas; no reescribe lo anterior. \\ \hline
\end{tabularx}
\caption{Revisiones del informe.}
\end{table}

\subsection{Cómo leer los identificadores}

Los requerimientos con prefijo \texttt{F-} y \texttt{NF-} son los del PRD de la Memoria Técnica,
y los \texttt{SR-} son los del SRD de la misma memoria. Los criterios \texttt{PoC-} son de este
informe: cada uno indica de cuál requerimiento deriva, o si es un criterio de diseño nuevo.
